\section{放射量に基づく段階的最適化}
\label{sec:optimization}

\subsection{損失をどこで計算するか}
\label{sec:lossdomain}
再照明の測光損失は、LDR 撮影データの本来の形式である、表示用に符号化された画像上で計算されることが多い。
本研究では、HDR データの表示に純粋なべき乗則 $\Gamma(L)=L^{1/\gamma}$、$\gamma=2.2$ を用い、LDR データの線形化にはその逆変換を用いる。このとき、予測放射輝度 $\hat L$ に対する $\ell_1$ 損失の勾配の大きさは
\begin{equation}
  \Gamma'(\hat L)=\tfrac{1}{\gamma}\,\hat L^{\,1/\gamma-1},
  \label{eq:slope}
\end{equation}
となり、残差の大きさに依存しない。
したがって、白の一パーセントの値を予測した画素には、残差がごく小さくても、白を予測した画素の十二倍の勾配が与えられる。
\cref{eq:radiance} の予測は学習する複数の因子の積であるため、学習初期に画素が暗いのは、可視性、法線、材質のいずれかがまだ誤っていることだけが原因の場合も多い。本研究では、そのような画素が幾何形状に到達する勾配を支配し、最適化器が細い構造を平滑化することでそれを解消する、という仮説を立てる。
\cref{sec:exp_optimization} の対照実験は、本研究のモデルにおいて、損失を計算する領域が細い幾何形状の形成可否を決めることを示す。ただし、その機序を切り分けたものではなく、$\Gamma$ 内へのオフセットの導入や損失の再重み付け~\cite{mildenhall2022nerf} といった対策も未検証である。

\subsection{三つの段階}
\label{sec:stages}
32 枚の学習用シルエットから視体積交差法で推定した形状の内部に、20k 個のガウシアンを一様に配置する。ある点がシルエットの少なくとも 90\% の内側に投影される場合にその点を採用し、法線はシーン中心から外向きに初期化する。

\paragraph{I．幾何形状のウォームアップ（8k 回の反復）。}
ネットワークを用いない点光源モデルで、各ガウシアンの色を $\softplus(\vect b_i)\,\bar E(\mu_i)\,\psi(\vect n_i\cdot\vect l_i)$ とし、幾何形状、ベースカラー、法線だけを最適化する。
レンダラの出力は線形のまま、目標値を変化させる。ウォームアップの最初の三分の二では表示用に符号化した画像 $\Gamma(Y)$ を目標とし、その後、べき指数を $1/\gamma$ から $1$ へ線形に変化させる。
これにより、幾何形状はダイナミックレンジを圧縮した目標値に対して形成され、連続的に線形放射輝度へ移行する。
\gsthree と SSD-GS は HDR 学習画像に類似の漸進的移行を適用する~\cite{bi2024gs3,zheng2026ssdgs}。本研究ではこれをすべてのデータ群に適用し、その後の同時学習も線形放射輝度上で行う。

\paragraph{II．同時学習（22k 回の反復）。}
\cref{eq:radiance} の完全モデルを、ガウシアン受光点と線形放射輝度上の損失を用いて学習する。高密度化と枝刈りは反復 25k まで 3DGS~\cite{kerbl20233d} に従って行い、ガウシアン数の上限を 400k とする。

\paragraph{III．外観の微調整（8k 回の反復）。}
幾何形状とカメラ校正を固定し、受光点を画素に切り替え、損失を画像の評価に用いる表示領域へ戻す。
幾何形状を固定した場合、本研究の実験では、暗い画素に対する重みの増加が両種の受光点で三つの評価指標すべてを改善する（\cref{sec:exp_optimization}）。

すべての段階で、現在の目標領域上で計算した $0.8\,\ell_1+0.2\,(1-\mathrm{SSIM})$ を最小化する。さらに、ベンチマークが提供する前景マスク $\alpha$ を用いたレンダリング不透明度の損失 $0.5\,\|\hat\alpha-\alpha\|_1$ と、コード $\vect f$ に対する弱い $\ell_2$ 正則化を加える。

\subsection{実装の詳細}
\label{sec:implementation}
\ours は PyTorch と gsplat~\cite{ye2025gsplat} を使用し、$512^2$ テクセル、$K=7$ 層、$C=8$ 光束チャネルのアトラスと、幅 128（可視性は 64）の MLP を用いる。学習が $L_{\mathrm{tr}}\approx 0$ および $V=V_0$ から始まるように初期化し、Adam を使用する（ネットワークの学習率は $10^{-3}$ から約四分の一まで減衰、ガウシアンは 3DGS と同じ設定）。
実撮影データでは、校正済みの中心を固定して各学習カメラの回転も微調整する。各ステップ後に、その微調整が吸収しうる全体の並進成分を除去し、再構成を校正済み座標系に保つ。テストカメラと光源は一切調整しない。
各ベンチマークのデータ規約だけを変え、全 18 シーンに一つの設定を共通で用いる。一回の完全な学習は NVIDIA RTX 6000 Ada 一枚で 17 から 28 分（平均 21 分）を要する。
